\begingroup\fontsize{9.3}{10.7}\selectfont\setlength{\parskip}{1pt}
\section{Notation for the gravitational realizations}
The notation distinguishes the objects of the construction, the response
operators, and their physical realizations. The positive inner products
of the extensions and the Lorentzian metric have different roles.
\renewcommand{\arraystretch}{1.0}\begin{longtable}{@{}p{.23\linewidth}p{.72\linewidth}@{}}
\toprule
Symbol & Object and role\\ \midrule
\endhead
\(G_a,\hbar_a\) & Reciprocal sections of gravitation and action.
Their product determines the invariant area
\(\ell_P^2=\hbar_aG_a/c^3=L_\alpha^2\) in the radial realization.\\
\(L_*,L_\alpha=L\) & Longitudinal reference section and constructed
length \(L_\alpha=(5759/23040)\alpha^{16}L_*\).\\
\(E_L,t_P,t_0\) & Return energy \(E_L=\hbar_{\rm ret}c/L_\alpha\),
time \(t_P=L_\alpha/c\), and step \(t_0=\pi_{\rm HMT}t_P/54\).\\
\(R_X,\Lambda_X\) & Positive radial reader and conjugate length of
the same character; \(R_X\Lambda_X=L_\alpha^2I\).\\
\(R_h,\zeta_h\) & Areal horizon radius and coordinate
\(R_h/L_\alpha=\zeta_h>0\). They retain the horizon's geometric condition.\\
\(\chi_{\rm rad}=L_\alpha/E_L\) & Dimensional ratio between radius
and energy. Its units and role differ from those of record \(K\)
and of \(\kappa=8\pi G/c^4\).\\
\(\gamma_{\rm HMT}\) & Evaluation of the hexad--octad functional,
reused as a coefficient of the Holst operator.\\
\(k_B,\Theta\) & Information-conversion coefficient and temperature
in the stated thermal section.\\
\(P_5,\Gamma_5,\Lambda_n\) & Projection onto the five coordinates,
tensor isometry, and rank-two transverse reduction.\\
\(c_s,c_g\) & Translational-memory cocycle and return component
of oriented transport.\\
\(\beta_m^\square,\Lambda_m^\square\) & Screen soldering and
energetic response, with their quadratic identity and scale.\\
\(\mathsf C,\mathsf L,f_{\min}\) & Transported incidence,
Gramian, and minimum-energy extension for the boundary datum.\\
\(s,\sigma,\mathsf M_{\rm cel}\) & Covector of the differential, material
current, and pairing that converts their coordinates.\\
\(e,\omega,\mathring\omega,K,T\) & Co-tetrad, connection,
Levi--Civita connection, contorsion, and torsion.\\
\(\mathsf P_\gamma,\mathsf L_e\) & Internal Holst operator and
linear map from torsion to the connection equation.\\
\(Q,F,\Delta S\) & Quadratic gravitational term, material action,
and effective correction after eliminating contorsion.\\
\(\mathsf H\) & Hessian of the composed action controlling the
local continuation of a torsional branch.\\
\(M_n,a,a_{\rm ref}\) & Memory, scale factor, and reference
normalization of its homogeneous realization.\\
\(s_0,\rho_0,w\) & Torsional amplitude, reference material density,
and barotropic parameter of the bounce domain.\\
\(\varrho,V_c,\widehat{\mathscr Q}_{\rm sp}\) & Normalized material state, comoving
volume, and normally ordered quartic observable; they determine
the functional \(s_0[\varrho,V_c]\).\\
\(C_\gamma\) & Coefficient of the effective spinorial interaction
obtained by eliminating contorsion.\\
\(\tau,t_{\rm b}\) & Characteristic duration and time of the minimum
of the homogeneous dust solution.\\
\(F_0,R,a_{\rm b}\) & Friedmann function, perturbation, and
transverse threshold of homogeneous evolution.\\
\(T_\Lambda,T_0\) & Trace and traceless parts of the residual tensor;
their divergences record exchange between sectors.\\
\(\mathsf{G}_{\rm ret},\mathsf{G}_{\rm adv}\) & Retarded and advanced propagators.
Here the letter \(G\) denotes operators, not the gravitational constant.\\
\(P_{\rm abs}\) & Projection onto the sources satisfying equality
of the two causal readings at the boundary.\\
\bottomrule
\end{longtable}

\par\endgroup
